% !TeX encoding = UTF-8
% Fuente editable del cuadernillo Ecoaldea Sostenible.
\documentclass{../plantilla/hvtbook}

\HVTTitulo{Ecoaldea Sostenible}
\HVTSubtitulo{Documento técnico para articular comunidad, territorio, vivienda, agua, saneamiento, energía, materiales, residuos, alimentación, movilidad, biodiversidad y gobernanza en un hábitat medible y resiliente.}
\HVTNivel{Proyecto de ecoaldea en Boyacá}
\HVTSerie{Publicación técnica}
\HVTNumero{03}
\HVTAccion{Explorar capítulos}
\HVTDescripcion{Publicación técnica del proyecto de investigación y diseño de una ecoaldea sostenible en Boyacá, desarrollado por Hidrógeno Verde Turquesa.}
\HVTCanonical{https://hidrogenoverdeturquesa.com/proyecto-ecoaldea}
\HVTRegreso{Volver a proyectos}{/\#portfolio}
\HVTPortada{../../images/portfolio/paneles.jpg}{Figura 1. Hábitat comunitario entendido como un sistema territorial integrado.}

\begin{document}
\HVTMakeTitle

\begin{HVTPage}{Contenido}{Ruta del proyecto}{Un territorio, no una suma de casas}
\HVTLead{La ecoaldea se diseña desde las personas y la capacidad del lugar. Cada subsistema debe funcionar por separado y, al mismo tiempo, relacionarse de forma segura con los demás.}
\begin{HVTTaskOrdered}{Capítulos}
  \item Comunidad y propósito.
  \item Diagnóstico territorial.
  \item Plan maestro y ocupación.
  \item Diseño bioclimático.
  \item Abastecimiento y balance de agua.
  \item Saneamiento y aguas residuales.
  \item Energía y servicios esenciales.
  \item Materiales y construcción.
  \item Residuos y economía circular.
  \item Agroecología y alimentación.
  \item Movilidad, seguridad y resiliencia.
  \item Gobernanza, operación y seguimiento.
\end{HVTTaskOrdered}
\end{HVTPage}

\begin{HVTPage}{Capítulo 1}{Comunidad}{Propósito y participación}
\HVTLead{El proyecto comienza identificando habitantes, propietarios, vecinos, autoridades y organizaciones relacionadas. La participación abierta y continua contribuye a la aceptación y a la sostenibilidad social del proyecto. \HVTCite{1}}
\begin{HVTFlow}
  \HVTFlowItem{1}{Escuchar}{Necesidades de vivienda, trabajo, cuidado, educación y vida comunitaria.}
  \HVTFlowItem{2}{Acordar}{Propiedad, aportes, derechos, responsabilidades y toma de decisiones.}
  \HVTFlowItem{3}{Documentar}{Actas, criterios de ingreso, manejo de conflictos y rendición de cuentas.}
\end{HVTFlow}
\begin{HVTTask}{Entregables}
  \item Mapa de actores y población prevista.
  \item Declaración de propósito y principios.
  \item Necesidades diferenciadas por grupo.
  \item Modelo preliminar de tenencia y uso.
  \item Mecanismo de consultas, quejas y respuesta.
\end{HVTTask}
\begin{HVTWarning}{Puerta de decisión 1}
No se define el diseño físico antes de acordar quiénes habitarán, qué necesitan y cómo se administrarán los bienes comunes.
\end{HVTWarning}
\end{HVTPage}

\begin{HVTPage}{Capítulo 2}{Territorio}{Diagnóstico del lugar}
\HVTLead{El terreno se evalúa como un sistema ambiental, social, jurídico y de infraestructura. El IDEAM aporta información oficial de clima, precipitación, temperatura, humedad, viento y radiación para una primera caracterización. \HVTCite{3}}
\begin{HVTTable}{Componente}{Información mínima}{Producto}
  \HVTTableRow{Jurídico}{Propiedad, linderos, servidumbres, uso del suelo y restricciones}{Concepto de viabilidad}
  \HVTTableRow{Físico}{Topografía, geología, geotecnia, pendientes y drenaje}{Mapa de aptitud}
  \HVTTableRow{Agua}{Fuentes, calidad, disponibilidad, rondas y escorrentía}{Balance preliminar}
  \HVTTableRow{Ecosistemas}{Coberturas, conectividad, especies y áreas sensibles}{Zonas de protección}
  \HVTTableRow{Riesgos}{Inundación, movimiento en masa, incendio, sismo y clima futuro}{Mapa multiamenaza}
  \HVTTableRow{Servicios}{Acceso, energía, comunicaciones, salud, educación y mercados}{Brechas y alternativas}
\end{HVTTable}
\end{HVTPage}

\begin{HVTPage}{Capítulo 2}{Territorio}{Condiciones de viabilidad}
\HVTText{Se deben consultar el ordenamiento territorial municipal, las autoridades ambientales y las licencias o permisos aplicables antes de comprometer el uso del predio.}
\begin{HVTWarning}{Puerta de decisión 2}
Excluir áreas de protección, amenaza no mitigable, conflicto jurídico o insuficiencia crítica de agua antes de elaborar el plan maestro.
\end{HVTWarning}
\end{HVTPage}

\begin{HVTPage}{Capítulo 3}{Plan maestro}{Ordenar la ocupación}
\HVTLead{Las directrices de ONU-Hábitat promueven territorios integrados, conectados, inclusivos y resilientes. El plan maestro traduce el diagnóstico en una estructura espacial verificable. \HVTCite{2}}
\begin{HVTTask}{Zonificación funcional}
  \item Conservación, restauración y corredores ecológicos.
  \item Vivienda y equipamientos comunitarios.
  \item Producción agroecológica y manejo de animales, si aplica.
  \item Captación, tratamiento y almacenamiento de agua.
  \item Saneamiento y aprovechamiento de residuos.
  \item Generación y almacenamiento energético.
  \item Movilidad peatonal, accesibilidad, emergencias y logística.
  \item Áreas de crecimiento sin ocupar zonas estratégicas.
\end{HVTTask}
\end{HVTPage}

\begin{HVTPage}{Capítulo 3}{Plan maestro}{Indicadores de ocupación}
\begin{HVTTable}{Indicador}{Unidad}{Decisión que apoya}
  \HVTTableRow{Población de diseño}{personas}{Capacidad de todos los servicios}
  \HVTTableRow{Densidad neta y bruta}{personas/ha}{Ocupación y costos de redes}
  \HVTTableRow{Área impermeabilizada}{\%}{Escorrentía y drenaje}
  \HVTTableRow{Distancia a servicios}{m o min}{Accesibilidad cotidiana}
  \HVTTableRow{Área protegida o restaurada}{m² o \%}{Integridad ecológica}
\end{HVTTable}
\begin{HVTWarning}{Puerta de decisión 3}
El plan maestro debe demostrar acceso seguro, capacidad de servicios y compatibilidad entre usos antes de diseñar las edificaciones.
\end{HVTWarning}
\end{HVTPage}

\begin{HVTPage}{Capítulo 4}{Bioclimática}{Confort con menor demanda}
\HVTLead{Primero se aplican suficiencia y estrategias pasivas; después eficiencia y energías renovables. El IPCC reúne estas estrategias para reducir demanda y emisiones en edificaciones. \HVTCite{5}}
\begin{HVTFlow}
  \HVTFlowItem{1}{Clima}{Temperatura, humedad, radiación, viento, lluvia y oscilación diaria.}
  \HVTFlowItem{2}{Forma}{Orientación, compacidad, patios, sombra y relación con el terreno.}
  \HVTFlowItem{3}{Envolvente}{Cubierta, muros, ventanas, masa térmica, aislamiento y hermeticidad.}
\end{HVTFlow}
\HVTEquation{\dot{Q}_{\text{trans}}=U\,A\left(T_{\text{ext}}-T_{\text{int}}\right)}{}
\HVTText{Esta relación estima la transferencia térmica estacionaria por un cerramiento. No sustituye la simulación dinámica, las ganancias solares, la ventilación, la humedad ni las cargas internas.}
\end{HVTPage}

\begin{HVTPage}{Capítulo 4}{Bioclimática}{Verificaciones}
\begin{HVTTask}{Verificaciones}
  \item Horas dentro del rango de confort definido.
  \item Iluminación natural y control de deslumbramiento.
  \item Ventilación y calidad del aire interior.
  \item Riesgo de condensación y humedad.
  \item Ruido, privacidad y accesibilidad universal.
  \item Cumplimiento de la Resolución 0194 de 2025 y demás requisitos aplicables a construcción sostenible en Colombia. \HVTCite{4}
\end{HVTTask}
\end{HVTPage}

\begin{HVTPage}{Capítulo 5}{Agua}{Abastecimiento y balance hídrico}
\HVTLead{La seguridad hídrica se diseña desde la disponibilidad real, la calidad requerida para cada uso y la variabilidad estacional, no únicamente desde el tamaño de un tanque.}
\begin{HVTTable}{Flujo}{Datos}{Unidad típica}
  \HVTTableRow{Demanda}{Consumo humano, cocina, higiene, limpieza y riego}{L/persona·día o m³/día}
  \HVTTableRow{Oferta}{Acueducto, fuente autorizada, lluvia y reúso validado}{m³/día}
  \HVTTableRow{Calidad}{Parámetros físicos, químicos y microbiológicos}{según uso}
  \HVTTableRow{Almacenamiento}{Autonomía, estacionalidad, incendio y pérdidas}{m³}
\end{HVTTable}
\HVTEquation{V_{\text{captable}}=P\,A\,C}{}
\HVTText{Si $P$ se expresa en metros, $A$ en m² y $C$ es un coeficiente adimensional, el resultado es m³ de agua de lluvia potencialmente captada antes de descontar primeras aguas, rebose y otras pérdidas.}
\end{HVTPage}

\begin{HVTPage}{Capítulo 5}{Agua}{Balance de almacenamiento}
\HVTEquation{V_{t+1}=V_t+V_{\text{entrada}}-V_{\text{demanda}}-V_{\text{pérdidas}}}{}
\begin{HVTWarning}{Precisión sanitaria}
Captar agua no garantiza su potabilidad. La fuente, tratamiento, almacenamiento y vigilancia deben diseñarse según la calidad requerida y la normativa colombiana aplicable.
\end{HVTWarning}
\end{HVTPage}

\begin{HVTPage}{Capítulo 6}{Saneamiento}{Barrera sanitaria completa}
\HVTLead{La OMS recomienda gestionar el saneamiento como una cadena segura desde el inodoro o punto de generación hasta el tratamiento, disposición o reúso final, controlando la exposición humana. \HVTCite{7}}
\begin{HVTFlow}
  \HVTFlowItem{1}{Separar}{Aguas lluvias, grises, negras y flujos con grasas o químicos.}
  \HVTFlowItem{2}{Tratar}{Definir procesos según caudal, carga, suelo, clima y objetivo de vertimiento o reúso.}
  \HVTFlowItem{3}{Vigilar}{Muestreo, operación, lodos, olores, vectores, contingencias y registros.}
\end{HVTFlow}
\HVTText{En Colombia, los sistemas de agua potable y saneamiento deben considerar el Reglamento Técnico RAS y sus modificaciones, además de permisos ambientales, criterios de vertimiento y condiciones locales. \HVTCite{6}}
\end{HVTPage}

\begin{HVTPage}{Capítulo 6}{Saneamiento}{Entregables y aprobación}
\begin{HVTTask}{Entregables}
  \item Caracterización y caudales de diseño.
  \item Alternativas centralizadas y descentralizadas.
  \item Balance de lodos y residuos del tratamiento.
  \item Evaluación de reúso con barreras sanitarias.
  \item Manual de operación, mantenimiento y emergencia.
\end{HVTTask}
\begin{HVTWarning}{Puerta de decisión 4}
No se aprueba la ocupación prevista sin demostrar abastecimiento seguro, tratamiento, disposición final y capacidad real de operación.
\end{HVTWarning}
\end{HVTPage}

\begin{HVTPage}{Capítulo 7}{Energía}{Servicios esenciales y microred}
\HVTLead{El sistema energético comienza reduciendo demanda y separando cargas críticas de cargas flexibles. La generación se dimensiona después con perfiles horarios y recursos medidos.}
\HVTEquation{E_{\text{demanda}}=\sum_i P_i\,t_i}{}
\HVTEquation{E_{\text{balance}}=E_{\text{generación}}+E_{\text{red}}-E_{\text{consumo}}-E_{\text{pérdidas}}}{}
\begin{HVTTable}{Nivel}{Ejemplos}{Criterio}
  \HVTTableRow{Crítico}{Agua, comunicaciones, seguridad, refrigeración esencial}{Continuidad y respaldo}
  \HVTTableRow{Prioritario}{Iluminación, cocina y equipos comunitarios}{Eficiencia y programación}
  \HVTTableRow{Flexible}{Bombeo, carga de vehículos y procesos productivos}{Desplazar hacia horas favorables}
\end{HVTTable}
\end{HVTPage}

\begin{HVTPage}{Capítulo 7}{Energía}{Estudios}
\begin{HVTTask}{Estudios}
  \item Curvas de carga y crecimiento.
  \item Recurso solar, eólico, biomasa u otras fuentes locales.
  \item Arquitectura conectada, aislada o híbrida.
  \item Almacenamiento, protecciones, puesta a tierra y calidad de energía.
  \item Operación en contingencia y recuperación.
  \item Costos de ciclo de vida y reposición.
\end{HVTTask}
\end{HVTPage}

\begin{HVTPage}{Capítulo 8}{Construcción}{Materiales seguros y durables}
\HVTLead{Un material no es sostenible solo por ser local o natural. Debe evaluarse su desempeño estructural, durabilidad, humedad, fuego, toxicidad, mantenimiento, disponibilidad y ciclo de vida.}
\begin{HVTTask}{Criterios de selección}
  \item Cumplimiento estructural, contra incendio y de salubridad.
  \item Energía y carbono incorporados.
  \item Extracción responsable y distancia de transporte.
  \item Contenido recuperado, reparación y desmontaje.
  \item Resistencia al clima, plagas, humedad y corrosión.
  \item Capacidad técnica de mano de obra y mantenimiento local.
  \item Gestión de residuos de construcción y demolición.
\end{HVTTask}
\HVTText{El análisis debe abarcar estructura, envolvente, acabados e instalaciones. Reducir la demanda operativa puede aumentar la importancia relativa de los impactos incorporados en los materiales, por lo que ambos deben revisarse. \HVTCite{5}}
\end{HVTPage}

\begin{HVTPage}{Capítulo 9}{Circularidad}{Prevenir antes de aprovechar}
\HVTLead{La Estrategia Nacional de Economía Circular de Colombia prioriza flujos de materiales, agua, energía, biomasa y construcción. La ecoaldea debe medir estos flujos y conservar su valor el mayor tiempo posible. \HVTCite{8}}
\begin{HVTFlow}
  \HVTFlowItem{1}{Evitar}{Compras innecesarias, desechables y materiales problemáticos.}
  \HVTFlowItem{2}{Reutilizar}{Reparación, intercambio, retorno y aprovechamiento interno seguro.}
  \HVTFlowItem{3}{Gestionar}{Separación, almacenamiento y entrega a operadores autorizados.}
\end{HVTFlow}
\HVTEquation{T_{\text{aprovechamiento}}=\frac{m_{\text{aprovechada}}}{m_{\text{generada}}}\,100\,\text{por ciento}}{}
\HVTText{La tasa debe informar qué materiales incluye, cómo se pesan y qué destino verificable reciben. Los residuos peligrosos, sanitarios, electrónicos y especiales requieren rutas diferenciadas.}
\end{HVTPage}

\begin{HVTPage}{Capítulo 9}{Circularidad}{Infraestructura}
\begin{HVTTask}{Infraestructura}
  \item Puntos de separación claramente identificados.
  \item Área seca para reciclables.
  \item Manejo controlado de orgánicos.
  \item Almacenamiento temporal seguro de residuos especiales.
  \item Registro de generación, aprovechamiento y disposición.
\end{HVTTask}
\end{HVTPage}

\begin{HVTPage}{Capítulo 10}{Agroecología}{Alimentación y paisaje productivo}
\HVTLead{La FAO propone diversidad, sinergias, reciclaje, eficiencia, resiliencia, conocimiento compartido, valores sociales, cultura alimentaria, gobernanza y economía solidaria como elementos interdependientes de la agroecología. \HVTCite{9}}
\begin{HVTTask}{Diseño productivo}
  \item Diagnóstico de suelo, agua, clima y conocimiento local.
  \item Dietas, demanda alimentaria y capacidades de trabajo.
  \item Diversidad de cultivos, rotaciones y especies adaptadas.
  \item Agroforestería, polinizadores y conectividad ecológica.
  \item Compostaje y bioinsumos con control de calidad.
  \item Riego eficiente y balance hídrico estacional.
  \item Manejo integrado de plagas y bioseguridad.
  \item Transformación, almacenamiento, intercambio y comercialización.
\end{HVTTask}
\end{HVTPage}

\begin{HVTPage}{Capítulo 10}{Agroecología}{Alcance productivo}
\begin{HVTWarning}{Alcance}
Producción local no significa autosuficiencia automática. Las metas deben derivarse de área útil, agua disponible, productividad medida, dieta, estacionalidad, trabajo y riesgos.
\end{HVTWarning}
\end{HVTPage}

\begin{HVTPage}{Capítulo 11}{Resiliencia}{Movilidad, seguridad y emergencias}
\HVTLead{El diseño debe permitir la vida cotidiana y la respuesta a emergencias para personas con diferentes edades y capacidades.}
\begin{HVTTable}{Sistema}{Verificación}{Contingencia}
  \HVTTableRow{Movilidad}{Rutas peatonales, accesibilidad, bicicletas, vehículos y logística}{Acceso de ambulancia y bomberos}
  \HVTTableRow{Incendio}{Vegetación, materiales, energía, agua y separación}{Detección, evacuación y control}
  \HVTTableRow{Agua}{Sequía, contaminación, inundación y falla de bombeo}{Reserva y fuente alternativa}
  \HVTTableRow{Energía}{Falla de red, batería, inversor o combustible}{Modo seguro para cargas críticas}
  \HVTTableRow{Salud}{Agua, alimentos, vectores, aire y atención}{Comunicación y traslado}
  \HVTTableRow{Comunidad}{Roles, visitantes y población vulnerable}{Plan de ayuda mutua}
\end{HVTTable}
\end{HVTPage}

\begin{HVTPage}{Capítulo 11}{Resiliencia}{Entregables}
\begin{HVTTask}{Entregables}
  \item Matriz de riesgos y responsables.
  \item Planos de evacuación y puntos de encuentro.
  \item Inventario de recursos y contactos.
  \item Protocolos de alarma, comunicación y recuperación.
  \item Programa de simulacros y revisión posterior.
\end{HVTTask}
\end{HVTPage}

\begin{HVTPage}{Capítulo 12}{Operación}{Gobernar, mantener y aprender}
\HVTLead{La sostenibilidad se demuestra durante la operación. Cada infraestructura necesita responsables, presupuesto, repuestos, registros y mecanismos de mejora.}
\begin{HVTTable}{Dimensión}{Indicador inicial}{Frecuencia}
  \HVTTableRow{Social}{Participación, acuerdos atendidos y satisfacción}{trimestral o anual}
  \HVTTableRow{Agua}{Consumo, calidad, fugas, almacenamiento y reúso}{continua y periódica}
  \HVTTableRow{Energía}{Demanda, generación, disponibilidad y pérdidas}{horaria y mensual}
  \HVTTableRow{Residuos}{Generación, aprovechamiento y destino}{mensual}
  \HVTTableRow{Alimentos}{Producción, agua, pérdidas y diversidad}{por ciclo}
  \HVTTableRow{Ecosistemas}{Cobertura, suelo, agua y especies indicadoras}{estacional o anual}
  \HVTTableRow{Economía}{Costos, fondos de reposición e ingresos}{mensual y anual}
\end{HVTTable}
\end{HVTPage}

\begin{HVTPage}{Capítulo 12}{Operación}{Documentos operativos}
\begin{HVTTask}{Documentos operativos}
  \item Reglamento de bienes comunes.
  \item Manuales por sistema.
  \item Plan de mantenimiento preventivo.
  \item Presupuesto de operación y reposición.
  \item Tablero de indicadores y trazabilidad de datos.
  \item Auditoría, aprendizaje y actualización del plan maestro.
\end{HVTTask}
\HVTHeading{Principio del proyecto}
\HVTText{Medir antes de decidir, integrar sin transferir riesgos y mantener lo que se construye.}
\end{HVTPage}

\begin{HVTPage}{Fuentes}{Referencias técnicas}{Documentos de consulta}
\begin{HVTReferences}
  \HVTReference{Banco Mundial. \emph{Environmental and Social Standard 10: Stakeholder Engagement and Information Disclosure}}{https://www.worldbank.org/en/projects-operations/environmental-and-social-framework/brief/environmental-and-social-standards}
  \HVTReference{ONU-Hábitat (2015). \emph{International Guidelines on Urban and Territorial Planning}}{https://unhabitat.org/international-guidelines-on-urban-and-territorial-planning}
  \HVTReference{IDEAM. \emph{Atlas Climatológico de Colombia y servicios meteorológicos}}{https://ideam.gov.co/nuestra-entidad/meteorologia}
\end{HVTReferences}
\end{HVTPage}

\begin{HVTPage}{Fuentes}{Referencias técnicas}{Documentos de consulta}
\begin{HVTReferences}
  \HVTReference{Ministerio de Vivienda, Ciudad y Territorio. \emph{Construcción sostenible y Resolución 0194 de 2025}}{https://minvivienda.gov.co/construccion-sostenible-y-cambio-climatico/construccion-sostenible}
  \HVTReference{IPCC (2022). \emph{Climate Change 2022: Mitigation of Climate Change, Chapter 9 --- Buildings}}{https://www.ipcc.ch/report/ar6/wg3/chapter/chapter-9/}
  \HVTReference{Ministerio de Vivienda, Ciudad y Territorio. \emph{Resolución 0330 de 2017 --- Reglamento Técnico para el Sector de Agua Potable y Saneamiento Básico (RAS)}}{https://minvivienda.gov.co/normativa/resolucion-0330-2017-0}
\end{HVTReferences}
\end{HVTPage}

\begin{HVTPage}{Fuentes}{Referencias técnicas}{Documentos de consulta}
\begin{HVTReferences}
  \HVTReference{Organización Mundial de la Salud (2018). \emph{Guidelines on Sanitation and Health}}{https://www.who.int/publications/i/item/9789241514705}
  \HVTReference{Ministerio de Ambiente y Desarrollo Sostenible. \emph{Estrategia Nacional de Economía Circular}}{https://www.minambiente.gov.co/asuntos-ambientales-sectorial-y-urbana/estrategia-nacional-de-economia-circular/}
  \HVTReference{Organización de las Naciones Unidas para la Alimentación y la Agricultura. \emph{Los 10 elementos de la agroecología}}{https://www.fao.org/agroecology/overview/the-10-elements-of-agroecology/the-10-elements-of-agroecology/es}
\end{HVTReferences}
\HVTText{Esta publicación presenta información técnica autorizada para divulgación pública. Los resultados detallados, datos operativos, análisis económicos y demás información estratégica del proyecto son confidenciales. Sus balances estructuran la ingeniería conceptual; el diseño definitivo requiere levantamientos, estudios ambientales y sociales, topografía, geotecnia, hidrología, arquitectura, ingeniería, permisos, análisis de riesgos y validación con la comunidad.}
\end{HVTPage}

\begin{HVTPage}{Colaboración}{Conversemos}{¿Necesitas investigar, formular o evaluar un proyecto relacionado?}
\HVTLead{Conversemos sobre tu necesidad de investigación, diseño territorial o evaluación de una ecoaldea o comunidad sostenible.}
\HVTContact{Consultar proyecto de ecoaldea por WhatsApp}{https://wa.me/573209574884}
\end{HVTPage}

\end{document}
